% Partitiones Oratoriae --- headnote
% partitiones-oratoriae, 54 BC, Rome

\textit{Partitiones Oratoriae, ``The Subdivisions of Oratory,'' is
Cicero's most compact and systematic handbook of rhetorical theory,
cast as a catechism between Cicero and his son Marcus. The son puts the
questions; the father answers, running through the whole art in tight
question-and-answer. In the Latin the labels are abbreviated and the
son is designated by the family name (\textit{Cicero}, that is, Cicero
the younger) against his father (\textit{Pater}); this translation
names the questioning son Marcus and the answering father Cicero, so
the reader is never in doubt who is speaking.}

\textit{The work was written in a moment of literary leisure away from
Rome, most likely in the 50s BC --- the conventional date is 54, though
some would bring it down as late as 46. Unlike the expansive dramatic
dialogue of De Oratore, this is a schoolroom epitome: Cicero drilling
his teenage son in the divisions he had himself laid out for him in
Greek, now turned into Latin. Its method is the Academic one --- the
closing words claim the whole descends from ``that Academy of ours'' ---
and it is built around a threefold division announced at the outset: the
faculty of the orator (discovery, arrangement, expression, memory, and
delivery), the speech (introduction, narrative, proof, and peroration),
and the question (the indefinite, which Cicero calls a general inquiry,
and the definite, which he calls a case --- the thesis and the
hypothesis of the Greek schools).}

\textit{The body then works through each division in order: the topics
from which arguments are drawn and the means of producing conviction and
stirring the emotions; the three kinds of case --- judicial,
deliberative, and panegyric; the long central treatment of arrangement
and the parts of a speech; the deliberative analysis of goods and of the
honourable and the advantageous; and finally the kinds of legal dispute
--- conjecture about the fact, definition, the letter of a law set
against its intent, ambiguity, and conflicting laws. It ends with the
father commending the whole as no more than a set of signposts to the
springs of philosophy, and the son's eager reply that of all his
father's gifts he looks for none greater.}
